% Nature-Methods-style report, generated by scripts/make_report.py (placeholders <<KEY>> are filled automatically)
\documentclass[9pt,twocolumn]{extarticle}
\usepackage[a4paper,margin=17mm,columnsep=6mm]{geometry}
\usepackage[scaled=0.95]{helvet}\renewcommand{\familydefault}{\sfdefault}
\usepackage[T1]{fontenc}
\usepackage{textcomp}
\usepackage{graphicx,booktabs,amsmath,amssymb,xcolor,siunitx,caption,titlesec,enumitem,microtype,tabularx,multirow}
\usepackage[hidelinks]{hyperref}
\usepackage[numbers,super,sort&compress]{natbib}
\definecolor{nmblue}{HTML}{0B3D91}
\captionsetup{font={footnotesize},labelfont={bf},labelsep={space},justification=justified,singlelinecheck=false}
\DeclareCaptionLabelSeparator{pipe}{ $|$ }
\captionsetup{labelsep=pipe}
\renewcommand{\figurename}{Fig.}
\titleformat{\section}{\normalfont\large\bfseries\color{nmblue}}{}{0pt}{}
\titleformat{\subsection}[runin]{\normalfont\bfseries}{}{0pt}{}[.]
\titlespacing*{\section}{0pt}{8pt}{3pt}
\titlespacing*{\subsection}{0pt}{4pt}{0.6em}
\setlength{\parindent}{1em}\setlength{\parskip}{0pt}
\renewcommand{\dbltopfraction}{0.95}\renewcommand{\dblfloatpagefraction}{0.6}\setcounter{dbltopnumber}{3}
\renewcommand{\topfraction}{0.95}\renewcommand{\textfraction}{0.05}\renewcommand{\floatpagefraction}{0.6}
\sisetup{detect-all,text-micro=\textmu,math-micro=\text{\textmu}}
\begin{document}
\twocolumn[{%
\begin{flushleft}
{\small\color{nmblue}\textbf{Technical report} \hfill <<DATE>>}\\[4pt]
{\LARGE\bfseries Pixel-level deformable registration of H\&E histology to en-face optical coherence tomography of skin-cancer specimens\par}
\vspace{6pt}
{\normalsize H\&E--OCT registration pipeline v1 \textperiodcentered{} \texttt{/home/adlz/atesfet/dev/registration/v1}\par}
\vspace{8pt}
\end{flushleft}
\noindent\textbf{Correlating optical coherence tomography (OCT) with histopathology requires mapping every OCT pixel to the corresponding location on the H\&E section. We present a fully automatic, GPU-accelerated pipeline that removes background from both modalities, normalises H\&E staining, determines orientation (including mirroring) and scale without calibration of the histology scanner, and deformably registers H\&E onto the OCT pixel grid. Across <<NSAMPLES>> specimens (<<SAMPLELIST>>), the median target registration error measured on <<POOLN>> independently segmented lumen landmarks fell from <<POOLAFFMED>>\,\textmu m after affine alignment to <<POOLDEFMED>>\,\textmu m after deformable registration (90th percentile <<POOLDEFP90>>\,\textmu m; <<POOLDEFLT2>>\% of landmarks within one OCT pixel of 2\,\textmu m), with a tissue-mask Dice of <<DICEMIN>>--<<DICEMAX>>.}
\vspace{10pt}
}]

\section{Main}
Optical coherence tomography (OCT)\cite{huang1991} images the surface of excised skin-cancer specimens without labels, but its interpretation rests on co-registered histopathology. Here, en-face OCT projections (mean of the five strongest depth slices of the tissue slab, in dB above the noise floor; \SI{2}{\micro\metre} pixels) are paired with 4$\times$ brightfield scans of the H\&E section cut from the same face. The two modalities differ in contrast mechanism (back-scattering versus absorption), in sampling depth (an approximately \SI{20}{\micro\metre} OCT slab versus a few-micrometre section) and in geometry: sectioning, floating and mounting stretch, compress, rotate and can mirror the section. Generic registration toolkits\cite{klein2010,borovec2020} require manual initialisation for such pairs and give no direct accuracy readout. We therefore built a dedicated pipeline (Fig.~\ref{fig:pipeline}) and an independent accuracy measure based on natural landmarks.

\section{Results}
\subsection{Preprocessing} H\&E background (glass, the cardboard holder, air bubbles, pen ink) was removed by combining optical density, saturation and hue criteria and keeping the single tissue section; Macenko stain normalisation\cite{macenko2009} mapped each slide to a common reference (Fig.~\ref{fig:pipeline}b). OCT background was detected by texture rather than intensity: background regions---noise floor, empty acquisition tiles and mounting medium---contain a high local fraction of near-noise pixels, whereas tissue does not. This was robust where intensity thresholding fails, for example when floating debris or a neighbouring tissue piece is brighter than the specimen. The shared $350\times350$-pixel tile pattern of the OCT mosaic was removed by an additive flat-field (Fig.~\ref{fig:pipeline}d).

\subsection{Registration} Orientation was found by a batched GPU search over 432 hypotheses (mirroring $\times$ 36 rotations $\times$ 6 scales) maximising tissue-mask overlap, after which the six best distinct hypotheses were re-ranked by image similarity using modality-independent neighbourhood descriptors (MIND)\cite{heinrich2012}. This disambiguates nearly round sections whose outlines alone do not determine rotation, and mirrored sections (mirrored samples: <<MIRRORED>>). A multi-scale affine fit followed by a five-level B-spline free-form deformation\cite{rueckert1999} (control spacing \SIrange{800}{50}{\micro\metre}, finest level on the native \SI{2}{\micro\metre} OCT grid) produced smooth, fold-free transformations (Table~\ref{tab:data}). The H\&E pixel size, unknown from the scanner file, was recovered from the fit (<<HEPX>>\,\textmu m per pixel). Overlays (Fig.~\ref{fig:overview}) and full-resolution regions of interest with complementary checkerboards (Fig.~\ref{fig:rois}) show that follicles, vessels, clefts and the epidermal boundary coincide.

\subsection{Registration accuracy} Because no manual landmarks exist, we measured target registration error (TRE)\cite{fitzpatrick2001} on lumens: vessels, gaps and clefts are low-scattering in OCT and unstained in H\&E, so both modalities show them as near-zero intensity. Lumens were segmented independently in each modality and paired by unambiguous mutual nearest neighbour (within \SI{20}{\micro\metre}, similar area, distance-ratio test). Pooled over all specimens, <<POOLN>> landmark pairs gave a median TRE of <<POOLAFFMED>>\,\textmu m (90th percentile <<POOLAFFP90>>\,\textmu m; <<POOLAFFLT2>>\% $<$\SI{2}{\micro\metre}) after affine alignment and \textbf{<<POOLDEFMED>>\,\textmu m} (90th percentile <<POOLDEFP90>>\,\textmu m; <<POOLDEFLT2>>\% $<$\SI{2}{\micro\metre}, <<POOLDEFLT4>>\% $<$\SI{4}{\micro\metre}) after deformable registration (Fig.~\ref{fig:accuracy}, Table~\ref{tab:acc}). A second, denser measure---the residual translation that would further align \SI{400}{\micro\metre} patches---had a median of <<RESDEFMED>>\,\textmu m after deformable registration, that is, about one OCT pixel. The landmark TRE is a conservative estimate: OCT integrates a thicker slab than the H\&E section, so lumen centroids genuinely differ between the two images, and residual mismatches in lumen-dense regions contribute a flat tail to the error distribution.

\begin{figure*}[t]
\centering\includegraphics[width=\textwidth]{figures/fig1_pipeline.png}
\caption{\textbf{Pipeline on specimen 12umA.} \textbf{a}, Raw 4$\times$ H\&E scan with holder, bubbles and ink. \textbf{b}, Tissue mask applied and Macenko-normalised. \textbf{c}, Raw OCT en-face top-$k$ projection (dB above noise floor). \textbf{d}, OCT after texture-based background and debris removal and tile flat-field correction. \textbf{e}, H\&E deformably registered onto the OCT grid, 50/50 blend. \textbf{f},\textbf{g}, Complementary checkerboards (\SI{500}{\micro\metre} tiles); together they show every pixel of both modalities. Scale bar, 2\,mm.}
\label{fig:pipeline}
\end{figure*}

\begin{figure*}[t]
\centering\includegraphics[width=\textwidth,height=0.82\textheight,keepaspectratio]{figures/fig2_overview.png}
\caption{\textbf{Whole-section registration for all specimens.} Columns: OCT (fixed); registered stain-normalised H\&E; checkerboard A; complementary checkerboard B; false-colour overlay (OCT magenta, H\&E optical density green; aligned structures appear white). Scale bars, 2\,mm.}
\label{fig:overview}
\end{figure*}

\begin{figure*}[t]
\centering\includegraphics[width=\textwidth,height=0.86\textheight,keepaspectratio]{figures/fig3_rois.png}
\caption{\textbf{Zoomed regions of interest at native OCT resolution (\SI{2}{\micro\metre} pixels, \SI{500}{\micro\metre} fields).} For each specimen, one structure-rich interior region and one epidermal-edge region are shown. Identical pixel coordinates are used in every panel; the yellow grid is \SI{100}{\micro\metre}. The edge overlay shows OCT edges (magenta), H\&E edges (green) and coincident edges (white). Scale bars, \SI{100}{\micro\metre}.}
\label{fig:rois}
\end{figure*}

\begin{figure*}[t]
\centering\includegraphics[width=\textwidth]{figures/fig4_accuracy.png}
\caption{\textbf{Registration accuracy.} \textbf{a}, Cumulative distribution of lumen-landmark target registration error per specimen after affine (dotted) and deformable (solid) registration; the dashed line marks one OCT pixel (\SI{2}{\micro\metre}). \textbf{b}, Median (bars) and 90th-percentile (ticks) landmark TRE. \textbf{c}, Median residual patch translation from MIND block matching (\SI{400}{\micro\metre} patches, $\pm$\SI{40}{\micro\metre} search).}
\label{fig:accuracy}
\end{figure*}

\begin{table*}[t]
\caption{\textbf{Registration accuracy per specimen.} Dice, tissue-mask overlap. TRE, lumen-landmark target registration error (\textmu m). Residual, MIND block-matching residual translation (\textmu m). Aff., affine; def., deformable.}
\label{tab:acc}
\centering\footnotesize
\begin{tabular}{lccrccccccccc}
\toprule
 & \multicolumn{2}{c}{Dice} & & \multicolumn{5}{c}{Landmark TRE} & \multicolumn{3}{c}{Residual} \\
\cmidrule(lr){2-3}\cmidrule(lr){5-9}\cmidrule(lr){10-12}
Specimen & aff. & def. & $n$ & aff. median & def. median & def. p90 & \% $<$2\,\textmu m & \% $<$4\,\textmu m & aff. median & def. median & def. p90 \\
\midrule
<<TABLE1>>
\bottomrule
\end{tabular}

\vspace{10pt}
\caption{\textbf{Data and transformation summary.} Image sizes are width $\times$ height in pixels. H\&E pixel size was estimated from the fitted affine transform (geometric mean of its two scale factors). Displacement is the non-rigid component on top of the affine transform. $J$, Jacobian determinant of the deformation (folding where $J\le0$).}
\label{tab:data}
\centering\footnotesize
\begin{tabular}{lcccccccccc}
\toprule
Specimen & H\&E size & OCT size & H\&E px (\textmu m) & mirrored & rotation ($^\circ$) & displ. median / max (\textmu m) & $J$ range & folding (\%) & runtime (min) \\
\midrule
<<TABLE2>>
\bottomrule
\end{tabular}
\end{table*}

\section{Discussion}
The pipeline aligns H\&E to OCT with a median landmark error of about 1.5 OCT pixels (<<POOLDEFMED>>\,\textmu m) and a median patch residual of about one pixel, with no manual input, and exports both modalities and a common tissue mask in one physical frame together with complete provenance metadata. Accuracy was lowest for sample\_12 (Table~\ref{tab:acc}), whose OCT contains saturated acquisition tiles, whose pale H\&E yields few lumens, and whose lumens are dilated in the section relative to OCT; in its structure-rich regions lumen outlines nevertheless coincide. The largest non-rigid displacements (up to about \SI{110}{\micro\metre}) are confined to a thin band at the tissue rim, where detached stratum corneum and edge tears in the section are folded back onto the OCT outline, and a \SI{700}{\micro\metre} pattern of 5--20\,\textmu m displacements in the interior corrects stitching offsets between OCT acquisition tiles. Otherwise, the remaining errors are concentrated where the section is torn or folded, at the cut edges of the OCT field of view, and in low-contrast dermis where neither modality offers distinctive structure. Three limitations should be noted. First, the landmark TRE mixes registration error with genuine anatomical differences between a thick optical slab and a thin physical section, so it overestimates the true error. Second, the H\&E pixel size is inferred rather than calibrated; if the scanner's calibration becomes available it can be fixed in the configuration. Third, specimen sample\_11 was excluded: its OCT field contains a second tissue piece and mounting medium in contact with the section, and it was set aside for separate handling.

\section{Methods}
\subsection{Data} Each specimen comprises a 4$\times$ H\&E RGB(A) TIFF without physical calibration and an OCT en-face \texttt{recon\_xy\_topk} TIFF (float32, dB above the depth-dependent noise floor, NaN outside the reconstructed footprint, \SI{2}{\micro\metre} pixels; 0.7\,mm acquisition patches of $350\times350$ pixels).

\subsection{H\&E preprocessing} The glass brightness $I_0$ was estimated per channel as the median of bright, unsaturated pixels. Optical density $\mathrm{OD}=-\ln(I/I_0)$ was computed at 1/8 resolution; tissue was defined as $\sum\mathrm{OD}>0.12$, HSV saturation $>0.10$ and hue outside the tan (holder) and blue (ink) ranges, followed by opening, closing (radius 48 full-resolution pixels), retention of components $\ge$25\% of the largest, and hole filling. Macenko normalisation\cite{macenko2009} estimated the stain matrix from up to $10^6$ tissue pixels (OD $>0.15$ in all channels, 1st/99th angle percentiles) and mapped concentrations to the standard reference (H $=[0.5626,0.7201,0.4062]$, E $=[0.2159,0.8012,0.5581]$, maximum concentrations 1.9705 and 1.0308). Background was set to white.

\subsection{OCT preprocessing} The local fraction of near-noise pixels, $f=G_{\sigma}\ast[\,v<4\,\mathrm{dB}\ \vee\ v=\mathrm{NaN}\,]$ with $\sigma=\SI{160}{\micro\metre}$, was thresholded ($f<0.35$) at 1/4 resolution. After a \SI{20}{\micro\metre} opening, retention of large components and hole filling, debris attached by bridges narrower than about \SI{480}{\micro\metre} was cut by an opening followed by a nearest-component (Voronoi) assignment, and touching lobes were split by a distance-transform watershed\cite{beucher1993} when their neck was narrower than 0.6 times the smaller lobe's inscribed radius. The tile flat-field was the Gaussian-smoothed median over fully-tissue tiles of (tile $-$ tile median), subtracted in dB. Intensities were clipped to the 0.5th--99.5th percentiles inside the mask and scaled to $[0,1]$.

\subsection{Transformation model} Let $\mathbf{x}$ be OCT coordinates in micrometres. The H\&E pixel position is $\mathbf{p}=\mathbf{A}_0\big[(\mathbf{I}+\mathbf{M})(\mathbf{x}+\mathbf{u}(\mathbf{x})-\mathbf{c})+\mathbf{c}+\mathbf{d}\big]+\mathbf{t}_0$, where $\mathbf{A}_0,\mathbf{t}_0$ encode the initial similarity transform (scale, rotation, optional mirroring), $\mathbf{M},\mathbf{d}$ the affine refinement about the image centre $\mathbf{c}$, and $\mathbf{u}=\sum_\ell \mathcal{B}(\mathbf{C}_\ell)$ a sum of cubically interpolated control grids $\mathbf{C}_\ell$. Images were sampled with bilinear (optimisation) or bicubic (export) interpolation in PyTorch\cite{paszke2019} on an NVIDIA A40.

\subsection{Similarity} The moving image was the summed optical density of the normalised H\&E, so lumens and background are near zero in both modalities. The loss was $\mathcal{L}=\overline{w\,\|\mathrm{MIND}(F)-\mathrm{MIND}(M\circ T)\|^2}+0.5\,(1-\mathrm{Dice})$, with 8-neighbourhood MIND\cite{heinrich2012} (Gaussian radius 2 and dilation 2 pixels for initialisation and affine fitting; radius 1 and dilation 1 for deformation), a weight $w$ covering both tissue masks, and soft mask Dice.

\subsection{Optimisation} (i) Initialisation: 432 hypotheses (mirroring $\{\pm1\}$ $\times$ rotations \ang{-180}--\ang{170} in \ang{10} steps $\times$ 6 log-spaced pixel sizes from 1.3 to 2.3\,\textmu m) were optimised jointly (rotation, scale, translation; Adam, 200 iterations) for soft mask Dice at \SI{32}{\micro\metre}; the six best distinct hypotheses (differing in mirroring or by more than \ang{20}) were refined for 100 affine iterations and ranked by $\mathcal{L}$. (ii) Affine: \SI{32}{}, \SI{16}{} and \SI{8}{\micro\metre} levels (300/200/150 iterations). (iii) Mask refinement: the OCT mask was intersected with the affine-registered H\&E mask dilated by \SI{100}{\micro\metre}, and the affine fit was repeated if more than 0.5\% was removed. (iv) Deformation: levels at \SI{16}{}, \SI{8}{}, \SI{4}{}, \SI{2}{} and \SI{2}{\micro\metre} with control spacing \SI{800}{}, \SI{400}{}, \SI{200}{}, \SI{100}{} and \SI{50}{\micro\metre} (450/400/300/250/250 iterations); each level adds a new control grid while earlier grids stay fixed, and the regulariser is $\lambda_\ell\overline{\|\nabla\mathbf{u}\|^2}+10\,\overline{\max(0,0.1-\det J)}$ with $\lambda_\ell=1,0.5,0.25,0.1,0.05$. These settings were chosen in a sweep on 12umA scored by landmark TRE.

\subsection{Accuracy assessment} Lumens were segmented as $G_{1}\ast I<0.1$ within the tissue mask eroded by \SI{80}{\micro\metre}, with areas of \SIrange{60}{12000}{\square\micro\metre}; pairs were mutual nearest neighbours within \SI{20}{\micro\metre} with area ratio $>0.3$ that were unambiguous in both directions (nearest/second-nearest distance ratio $<0.5$, after Lowe), which rejects chance pairings in lumen-dense tissue. Block matching computed, for \SI{400}{\micro\metre} patches (50\% overlap) at \SI{4}{\micro\metre} resolution, the integer shift within $\pm$\SI{40}{\micro\metre} that minimised the MIND sum of squared differences, refined to sub-pixel precision by a parabolic fit; patches with $<$90\% tissue, indistinct minima or minima on the search border were excluded. Because block matching uses the optimised similarity measure, it is reported as a consistency check, whereas the landmark TRE is independent of the optimiser.

\subsection{Outputs} For each specimen, \texttt{output/} contains lossless TIFFs on (i) the native OCT grid (\SI{2}{\micro\metre}) and (ii) an H\&E-native grid over the same physical frame: OCT (dB), registered H\&E (normalised and original colours), and OCT, H\&E and common tissue masks (common mask $=$ OCT $\wedge$ registered H\&E). Every TIFF is ImageJ-calibrated and embeds the complete metadata, which is also written to a single \texttt{<sample>\_metadata.json}: inputs with SHA-256 checksums, pixel sizes, coordinate frame and its conversion to OCT scanner millimetres, transform parameters, preprocessing parameters, stain matrices, accuracy metrics and software versions. The full transform (\texttt{transform\_full.npz}) and a dense per-pixel map are kept for mapping annotations in either direction.

\subsection{Code and data availability} {\raggedright Code: \texttt{/home/adlz/atesfet/dev/registration/v1} (conda environment \texttt{reg\_env}). Registered data and per-specimen quality-control pages: \texttt{/home/adlz/atesfet/dev/data\_12um\_registered/<sample>/}.\par}

\begingroup\footnotesize
\begin{thebibliography}{10}\setlength{\itemsep}{0pt}
\bibitem{huang1991} Huang, D. \emph{et al.} Optical coherence tomography. \emph{Science} \textbf{254}, 1178--1181 (1991).
\bibitem{klein2010} Klein, S., Staring, M., Murphy, K., Viergever, M. A. \& Pluim, J. P. W. elastix: a toolbox for intensity-based medical image registration. \emph{IEEE Trans. Med. Imaging} \textbf{29}, 196--205 (2010).
\bibitem{borovec2020} Borovec, J. \emph{et al.} ANHIR: automatic non-rigid histological image registration challenge. \emph{IEEE Trans. Med. Imaging} \textbf{39}, 3042--3052 (2020).
\bibitem{macenko2009} Macenko, M. \emph{et al.} A method for normalizing histology slides for quantitative analysis. In \emph{Proc. IEEE Int. Symp. Biomed. Imaging} 1107--1110 (2009).
\bibitem{heinrich2012} Heinrich, M. P. \emph{et al.} MIND: modality independent neighbourhood descriptor for multi-modal deformable registration. \emph{Med. Image Anal.} \textbf{16}, 1423--1435 (2012).
\bibitem{rueckert1999} Rueckert, D. \emph{et al.} Nonrigid registration using free-form deformations: application to breast MR images. \emph{IEEE Trans. Med. Imaging} \textbf{18}, 712--721 (1999).
\bibitem{fitzpatrick2001} Fitzpatrick, J. M. \& West, J. B. The distribution of target registration error in rigid-body point-based registration. \emph{IEEE Trans. Med. Imaging} \textbf{20}, 917--927 (2001).
\bibitem{beucher1993} Beucher, S. \& Meyer, F. The morphological approach to segmentation: the watershed transformation. In \emph{Mathematical Morphology in Image Processing} (ed. Dougherty, E. R.) 433--481 (Marcel Dekker, 1993).
\bibitem{paszke2019} Paszke, A. \emph{et al.} PyTorch: an imperative style, high-performance deep learning library. In \emph{Adv. Neural Inf. Process. Syst.} \textbf{32}, 8024--8035 (2019).
\end{thebibliography}
\endgroup
\end{document}
